\documentclass[11pt]{article}

\usepackage[T1]{fontenc}
\usepackage{lmodern}
\usepackage{amsmath,amssymb,amsthm}
\usepackage[margin=1in]{geometry}
\usepackage[hidelinks]{hyperref}

\newtheorem{theorem}{Theorem}
\newtheorem{remark}[theorem]{Remark}
\newcommand{\supp}{\operatorname{supp}}
\newcommand{\Frac}{\operatorname{Frac}}

\title{The Exact Six-Term Eliminant in a Two-State Affine-Linear Differential-Input System}
\author{Xiaojian Zeng\\
PBC School of Finance, Tsinghua University, 100190\\
\texttt{xiaojianzeng@foxmail.com}}
\date{}

\begin{document}

\maketitle

\begin{abstract}
For the generic dense two-state affine-linear system with one differential
input and an affine-linear output, we determine the input--output eliminant
over independent, derivation-constant coefficient indeterminates.  On an
explicit nonempty Zariski-open coefficient locus, its order is two and its
exact support is
\(\{1,u,u',y,y',y''\}\).  Thus, in the first differential-parameter row
reported by Jeronimo and Lanciano, the missing one of the seven support-bound
candidates is \(u''\), and the additional exact support equality is
\(\ell_2=0\).  The proof identifies the full defining ideal of the
order-two input--output image, rather than inferring genericity from
finite-field samples.  The result is confined to this affine-linear,
two-state, characteristic-zero row.
\end{abstract}

\section{Context and statement}

Input--output elimination replaces hidden states of a polynomial dynamical
system by differential relations among its measured inputs and outputs; this
is also a basic operation in structural-identifiability computations
\cite{DongEtAl2023}.  Recent support bounds make evaluation--interpolation
reconstruction possible first for coordinate projections and systems with
constant symbolic parameters \cite{MukhinaPogudin2025Projection,MukhinaPogudin2026},
and then for freely varying differential inputs \cite{JeronimoLanciano2026}.

Theorem~3.7 of \cite{JeronimoLanciano2026} gives a finite candidate support.
For two states, one input, and degree data
\[
 [D_x,D_u,D]=[1,1,1],\qquad [d_x,d_u,d]=[1,0,1],
\]
the first differential-input row of that paper reports seven candidates but
six stable terms in finite-field experiments.  We give the corresponding
coefficient-generic calculation and explain the discrepancy exactly.  Here
``generic'' always refers to an explicitly exhibited nonempty Zariski-open
set, not to a finite sample.  This same distinction between a generic
characteristic-zero statement and its exceptional algebraic locus is useful
in other exact polynomial computations; see, for example,
\cite{Zeng2026Bezout}.  That work concerns algebraic complexity rather than
differential elimination and is cited only for this methodological parallel.

Let
\[
 R=\mathbb{Q}[a_{11},a_{12},a_{21},a_{22},b_1,b_2,c_1,c_2,p_1,p_2,q],
 \qquad F=\Frac(R),
\]
where every generator of \(R\) is a derivation constant.  Set
\[
 A=\begin{pmatrix}a_{11}&a_{12}\\a_{21}&a_{22}\end{pmatrix},\quad
 b=\binom{b_1}{b_2},\quad c=\binom{c_1}{c_2},\quad
 p^T=(p_1,p_2),
\]
and consider
\begin{equation}\label{eq:system}
 x'=Ax+bu+c,\qquad y=p^Tx+q.
\end{equation}
Put
\[
 \tau=\operatorname{tr}(A),\qquad \delta=\det(A),\qquad
 \Omega=\det\!\begin{pmatrix}p^T\\p^TA\end{pmatrix}
\]
and
\begin{equation}\label{eq:P}
 P=\Omega\,\tau\,\delta\,(p^Tb)
 \bigl(p^T(\tau I-A)b\bigr)
 \bigl(p^T(\tau I-A)c-\delta q\bigr).
\end{equation}

\begin{theorem}\label{thm:main}
The polynomial \(P\) is nonzero.  On the nonempty Zariski-open locus
\(P\ne0\), the input--output eliminant of \eqref{eq:system} has order two
and, normalized to be monic in \(y''\), equals
\begin{equation}\label{eq:E}
 E=y''-\tau y'+\delta y
   +p^T(\tau I-A)b\,u-p^Tb\,u'
   +p^T(\tau I-A)c-\delta q.
\end{equation}
It is primitive and irreducible, and
\[
 \supp(E)=\{1,u,u',y,y',y''\}.
\]
The seven nonnegative integer points allowed by the specialized
Theorem~3.7 inequalities of \cite{JeronimoLanciano2026} correspond to
\(\{1,u,u',u'',y,y',y''\}\).  Hence the missing monomial is \(u''\), and
every exponent vector in \(\supp(E)\) satisfies
\[
 \ell_2=0,
 \qquad\text{equivalently}\qquad \deg_{u''}E=0.
\]
\end{theorem}

The quantifier in Theorem~\ref{thm:main} may be read over any algebraically
closed characteristic-zero extension field.  Its scope does not include
nonlinear or sparse systems, time-varying model coefficients, positive
characteristic, or additional states or inputs.

\section{Proof of the exact eliminant}

Differentiating the output after substituting the state equation gives
\begin{align}
 y'  &=p^TAx+p^Tb\,u+p^Tc,\label{eq:yp}\\
 y'' &=p^TA^2x+p^TAb\,u+p^TAc+p^Tb\,u'.\label{eq:ypp}
\end{align}
The two-by-two Cayley--Hamilton identity
\(A^2-\tau A+\delta I=0\) converts \eqref{eq:ypp} into
\eqref{eq:E}.  Thus \(E\) vanishes identically over the independent
coefficient ring; no specialization has been used.

We next prove that this identity is the defining eliminant and has minimal
order.  Direct expansion gives
\begin{equation}\label{eq:omega}
 \Omega=p_1^2a_{12}+p_1p_2(a_{22}-a_{11})-p_2^2a_{21},
\end{equation}
which is a nonzero element of \(R\).  Over
\(L=\Frac(F\{u\})\), equations \eqref{eq:system} and \eqref{eq:yp} yield
\begin{equation}\label{eq:obs}
 \binom{y-q}{y'-p^Tb\,u-p^Tc}
 =\begin{pmatrix}p^T\\p^TA\end{pmatrix}x.
\end{equation}
Since \(\Omega\ne0\) in \(F\), this is an invertible affine change of
coordinates over \(L\).  Consequently \(y,y'\) are algebraically independent
over \(L\), while \eqref{eq:E} makes \(y,y',y''\) dependent.  The generic
order is therefore exactly two.

For the defining-ideal assertion, set \(B=F[u_0,u_1,u_2]\) and let
\[
 \psi:B[Y_0,Y_1,Y_2]\longrightarrow B[x_1,x_2]
\]
send \(Y_i\) to the evaluated \(i\)-th output derivative and fix the
\(u_i\).  Equation \eqref{eq:obs} shows that the restriction
\(B[Y_0,Y_1]\to B[x_1,x_2]\) is an isomorphism.  Write
\[
 \alpha=p^T(\tau I-A)b,\qquad \beta=p^Tb,\qquad
 \gamma=p^T(\tau I-A)c.
\]
Under this isomorphism, \(Y_2\) is the affine-linear function
\[
 \tau Y_1-\delta Y_0-\alpha u_0+\beta u_1-\gamma+\delta q.
\]
Division in the monic variable \(Y_2\) proves that
\(\ker(\psi)=(E)\).  Hence the image is the closed graph cut out by
\(E\), not merely a subset on which \(E\) happens to vanish.  Its
\(Y_2\)-coefficient is one, so \(E\) is primitive; being monic and linear in
\(Y_2\), it is irreducible.  This is precisely the normalized eliminant in
the sense used by \cite{JeronimoLanciano2026}.

It remains to establish support equality, rather than only containment.  The
six coefficients, in the order
\[
 y'',\ y',\ y,\ u',\ u,\ 1,
\]
are
\begin{equation}\label{eq:coefficients}
 1,\quad-\tau,\quad\delta,\quad-p^Tb,\quad
 p^T(\tau I-A)b,\quad p^T(\tau I-A)c-\delta q.
\end{equation}
At the exact specialization
\[
 A=\begin{pmatrix}1&1\\1&2\end{pmatrix},\quad
 p^T=(1,0),\quad b=(1,1)^T,\quad c=(1,0)^T,\quad q=0,
\]
the factors in \eqref{eq:P}, including \(\Omega\), take the values
\[
 (\Omega,\tau,\delta,p^Tb,p^T(\tau I-A)b,
 p^T(\tau I-A)c-\delta q)=(1,3,1,1,1,2).
\]
Thus \(P\) is not the zero polynomial.  On \(P\ne0\), every coefficient in
\eqref{eq:coefficients} is nonzero, proving the asserted six-term support.

Finally, use the exponent vector
\[
 (\ell_0,\ell_1,\ell_2,e_0,e_1,e_2)
\]
for the variables \(u,u',u'',y,y',y''\), in that order.  The specialized
support inequalities are
\begin{align*}
 \ell_0+\ell_1+\ell_2+e_0+e_1+e_2&\le1,&
 e_0+e_1+e_2&\le1,\\
 \ell_0+\ell_1+\ell_2+e_1+2e_2&\le3,&
 \ell_1+e_1+2e_2&\le3,\\
 \ell_2+e_2&\le1.&&
\end{align*}
The first inequality leaves only the origin and the six unit vectors, and
each satisfies the remaining inequalities.  Hence the candidate monomials
are exactly \(1,u,u',u'',y,y',y''\).  Formula \eqref{eq:E} contains all but
\(u''\), completing the proof of Theorem~\ref{thm:main}. \qed

\begin{remark}
Corollary~3.11 of \cite{JeronimoLanciano2026} observes that, when the output
is a function of one distinguished state coordinate, the eliminant does not
depend on the highest input derivative.  On \(\Omega\ne0\), an invertible
linear state-coordinate change can make the affine output in
\eqref{eq:system} the distinguished coordinate.  Thus the absence of
\(u''\) is consistent with that general structural observation.  The
calculation above additionally supplies the normalized eliminant, proves its
minimality and defining-ideal property, and establishes generic nonvanishing
of all six remaining coefficients for the reported row.
\end{remark}

\section{Exact verification and limitations}

An accompanying SymPy script expands the Cayley--Hamilton relation and checks
the inverse-coordinate identity after multiplication by \(\Omega\).  It also
verifies the graph equation, enumerates the seven integer points, and
evaluates the open-set witness using exact arithmetic.  These checks are
reproducibility aids; the proof is the coefficient-ring argument above and
does not rest on numerical or finite-field agreement.

This row-wise calculation does not sharpen the general support polytope for
nonlinear systems with differential inputs.  In particular, it does not address
the broader problem, identified in \cite{JeronimoLanciano2026}, of obtaining
more accurate support descriptions in that setting.

\paragraph{Computational and AI assistance.}
SymPy was used for exact expansion, integer-point enumeration, and replay of
the displayed specialization.  OpenAI Codex assisted with the symbolic proof
workflow, source screening, and manuscript preparation.  The author reviewed
the mathematical argument and cited-source statements and takes
responsibility for the contents.

\bibliographystyle{plain}
\bibliography{references}

\end{document}
